\section{Ethics Considerations}
\label{sec:ethics}

This appendix describes the ethical posture of the work in the
categories USENIX Security policy names for the Ethics Considerations
section.

\subsection{Human and animal subjects}

No human-subjects research and no animal-subjects research were
conducted. No IRB review or animal-care protocol was required. The
empirical evaluation in Section~\ref{sec:evaluation} is over receipts
produced by a software admission kernel; the measured artifacts are
signed canonical-JSON records and verifier-side mutation outcomes, not
data derived from human participants. No personally identifying
information was collected, analyzed, or stored at any point in the
construction or evaluation.

\subsection{Responsible disclosure}

The sensor-grounded admission construction does not introduce new
vulnerabilities. The construction is a verifier-side admission
discipline that consumes existing trusted-execution-environment
attestation primitives whose security properties are inherited from
the underlying substrate. No new attacks against attestation
primitives, signing keys, sensing hardware, or cryptographic
constructions are proposed.

The TEE-compromise literature cited in Section~\ref{sec:limits}
describes attacks that were disclosed by their original authors,
assigned CVE identifiers, and (in most cases) mitigated by hardware
vendors through microcode updates or successor generations:
Plundervolt~\cite{murdockPlundervolt2020},
Foreshadow~\cite{vanbulckForeshadow2018},
Downfall~\cite{moghimiDownfall2024}, and
Half-Double~\cite{koglerHalfDouble2022} are each cited in their
published, peer-reviewed forms. The paper acknowledges that TEE
compromise reduces the construction to single-key signing for the
duration of a compromise window, and names this reduction as a known
structural boundary rather than as a new finding. No
responsible-disclosure timeline was triggered by this work, and no
embargo applies to any artifact in the submission.

\subsection{Dual-use considerations}

The construction's primary beneficiary is the verifier, that is, the
polity admitting receipts under a constitution. An attacker who
controls the agent producing a receipt body gains nothing from the
substrate attesting its own sensing posture: the attestation is
adversarial information from the attacker's perspective, since it
constrains rather than enables admission. Forging an attestation
requires compromising the TEE-rooted attestation signing key, which
is the threat-model boundary named in Section~\ref{sec:intro} and
discussed in Section~\ref{sec:limits}.

The marginal dual-use concern is observability: a verifier sees more
about the signing kernel's substrate state than under the
unstrengthened predicate, and an adversary positioned in the verifier
role inherits that observability. This concern is intrinsic to any
attested-system design that exposes substrate state to relying
parties, and is not unique to the sensor-grounded construction. The
construction does not introduce a novel attack surface, novel
side-channel, or novel cryptographic primitive whose misuse would
benefit an attacker.

\subsection{Conflicts of interest}

Following USENIX Security double-blind review policy, conflict-of-interest
declarations are deferred to the camera-ready version. The authors
will disclose any relationships requiring declaration under USENIX
policy at that stage. The work does not benefit from any commercial
relationship whose disclosure is foreclosed by the anonymous-review
posture.
